\chapter{Problem Definition}
\label{chap:problem-def}

Speech Enhancement stands as a pivotal domain within signal processing, dedicated to refining speech signals captured under adverse conditions marked by noise or degradation. The core objective is the enhancement of quality and intelligibility in these speech signals.

The overarching aim of speech enhancement is to elevate the clarity and understandability of speech, particularly in challenging environments.

Our specific focus lies in addressing the challenges associated with speech enhancement in the context of online communication. This is a crucial area, as the quality of speech in virtual interactions significantly impacts user experience and effective communication.

Speech enhancement finds widespread application across diverse domains, including but not limited to:

\begin{itemize}
\item Mobile Communication
\item Teleconferencing
\item Speech Recognition Systems
\item Hearing Aids
\end{itemize}

In many of these applications, a fundamental requirement for a speech enhancement system is its ability to operate in real time with minimal latency (online), preferably on cost-effective hardware integrated into communication devices.

Our proposed approach involves leveraging Deep Learning methods specifically tailored for Real-Time Speech Enhancement. It is noteworthy that real-time implementations, by nature, demand heightened efficiency compared to their non-real-time counterparts.

In our pursuit of advancing the current state-of-the-art, we aim to enhance an existing speech enhancement method. We plan to achieve this by incorporating a crucial element—utilizing a reference of \textbf{clear speech} provided by \textbf{the device user}. This reference will serve as a valuable asset, aiding the model in effectively extracting and enhancing the user's voice even in the presence of substantial ambient noise. This innovative approach is anticipated to contribute significantly to the improvement of speech quality in online communication scenarios.